\begin{figure}[tbp]
\centering
\begin{tikzpicture}[x=1.65cm,y=1.3cm,>=Stealth,font=\small]
  \fill[blue!7] (0.45,-0.35) rectangle (5.55,0);
  \fill[orange!12] (1,0) rectangle (3,0.35);
  \fill[orange!12] (3,0) rectangle (5,0.35);
  \draw[->,gray] (0,-0.98)--(0,1.9) node[above]{height};
  \draw[->,gray] (0,-0.95)--(5.95,-0.95) node[right]{time};
  \draw[gray,dotted] (0,0)--(5.6,0);
  \node[left] at (0,0) {$H$};
  \node[left] at (0,-0.3) {$H-\varepsilon$};
  \node[left] at (0,0.3) {$H+\varepsilon$};
  \draw[very thick] (0.45,-0.55)--(1,0)
    .. controls (1.35,1.2) and (1.15,1.55) .. (1.65,1.35)
    .. controls (2.05,1.1) and (1.85,0.6) .. (2.2,1.0)
    .. controls (2.6,1.2) and (2.55,0.25) .. (3,0)
    .. controls (3.3,0.15) and (3.2,1.0) .. (3.65,0.8)
    .. controls (4.0,0.55) and (4.0,1.65) .. (4.35,1.5)
    .. controls (4.7,1.2) and (4.72,0.55) .. (5,0)
    --(5.55,-0.55);
  \draw[blue!70!black,thick,<->] (0.7,-0.3)--(5.3,-0.3);
  \node[blue!70!black,fill=blue!7,inner sep=1pt] at (3,-0.3)
    {outer seam};
  \draw[orange!75!black,thick,<->] (1.1,0.3)--(2.65,0.3);
  \draw[orange!75!black,thick,<->] (3.25,0.3)--(4.85,0.3);
  \node[orange!75!black] at (1.9,1.85) {adjacent seam};
  \node[orange!75!black] at (4.15,1.85) {adjacent seam};
  \foreach \x/\i in {1/1,3/2,5/3}{
    \fill (\x,0) circle (1.8pt);
    \node[fill=white,inner sep=1pt] at (\x,-0.65) {$t_\i$};
  }
\end{tikzpicture}
\caption{A triple contact in one contour coordinate. The visits
$t_1<t_2<t_3$ at height $H$ represent the same surface point.
Overlaps below $H$ join the outer visits; overlaps above $H$ join
adjacent visits. Their directions determine the relative offsets and
identify the middle run, whose exact minimum supplies the cutoff $H$.
Omitted gaps stay separated from the retained bands by a positive
height margin. This diagram shows the incidences schematically.}
\label{tree:fig:triple-contact}
\end{figure}
